% Options for packages loaded elsewhere
\PassOptionsToPackage{unicode}{hyperref}
\PassOptionsToPackage{hyphens}{url}
%
\documentclass[
]{article}
\usepackage{amsmath,amssymb}
\usepackage{iftex}
\ifPDFTeX
  \usepackage[T1]{fontenc}
  \usepackage[utf8]{inputenc}
  \usepackage{textcomp} % provide euro and other symbols
\else % if luatex or xetex
  \usepackage{unicode-math} % this also loads fontspec
  \defaultfontfeatures{Scale=MatchLowercase}
  \defaultfontfeatures[\rmfamily]{Ligatures=TeX,Scale=1}
\fi
\usepackage{lmodern}
\ifPDFTeX\else
  % xetex/luatex font selection
\fi
% Use upquote if available, for straight quotes in verbatim environments
\IfFileExists{upquote.sty}{\usepackage{upquote}}{}
\IfFileExists{microtype.sty}{% use microtype if available
  \usepackage[]{microtype}
  \UseMicrotypeSet[protrusion]{basicmath} % disable protrusion for tt fonts
}{}
\makeatletter
\@ifundefined{KOMAClassName}{% if non-KOMA class
  \IfFileExists{parskip.sty}{%
    \usepackage{parskip}
  }{% else
    \setlength{\parindent}{0pt}
    \setlength{\parskip}{6pt plus 2pt minus 1pt}}
}{% if KOMA class
  \KOMAoptions{parskip=half}}
\makeatother
\usepackage{xcolor}
\usepackage[margin=1in]{geometry}
\usepackage{color}
\usepackage{fancyvrb}
\newcommand{\VerbBar}{|}
\newcommand{\VERB}{\Verb[commandchars=\\\{\}]}
\DefineVerbatimEnvironment{Highlighting}{Verbatim}{commandchars=\\\{\}}
% Add ',fontsize=\small' for more characters per line
\usepackage{framed}
\definecolor{shadecolor}{RGB}{248,248,248}
\newenvironment{Shaded}{\begin{snugshade}}{\end{snugshade}}
\newcommand{\AlertTok}[1]{\textcolor[rgb]{0.94,0.16,0.16}{#1}}
\newcommand{\AnnotationTok}[1]{\textcolor[rgb]{0.56,0.35,0.01}{\textbf{\textit{#1}}}}
\newcommand{\AttributeTok}[1]{\textcolor[rgb]{0.13,0.29,0.53}{#1}}
\newcommand{\BaseNTok}[1]{\textcolor[rgb]{0.00,0.00,0.81}{#1}}
\newcommand{\BuiltInTok}[1]{#1}
\newcommand{\CharTok}[1]{\textcolor[rgb]{0.31,0.60,0.02}{#1}}
\newcommand{\CommentTok}[1]{\textcolor[rgb]{0.56,0.35,0.01}{\textit{#1}}}
\newcommand{\CommentVarTok}[1]{\textcolor[rgb]{0.56,0.35,0.01}{\textbf{\textit{#1}}}}
\newcommand{\ConstantTok}[1]{\textcolor[rgb]{0.56,0.35,0.01}{#1}}
\newcommand{\ControlFlowTok}[1]{\textcolor[rgb]{0.13,0.29,0.53}{\textbf{#1}}}
\newcommand{\DataTypeTok}[1]{\textcolor[rgb]{0.13,0.29,0.53}{#1}}
\newcommand{\DecValTok}[1]{\textcolor[rgb]{0.00,0.00,0.81}{#1}}
\newcommand{\DocumentationTok}[1]{\textcolor[rgb]{0.56,0.35,0.01}{\textbf{\textit{#1}}}}
\newcommand{\ErrorTok}[1]{\textcolor[rgb]{0.64,0.00,0.00}{\textbf{#1}}}
\newcommand{\ExtensionTok}[1]{#1}
\newcommand{\FloatTok}[1]{\textcolor[rgb]{0.00,0.00,0.81}{#1}}
\newcommand{\FunctionTok}[1]{\textcolor[rgb]{0.13,0.29,0.53}{\textbf{#1}}}
\newcommand{\ImportTok}[1]{#1}
\newcommand{\InformationTok}[1]{\textcolor[rgb]{0.56,0.35,0.01}{\textbf{\textit{#1}}}}
\newcommand{\KeywordTok}[1]{\textcolor[rgb]{0.13,0.29,0.53}{\textbf{#1}}}
\newcommand{\NormalTok}[1]{#1}
\newcommand{\OperatorTok}[1]{\textcolor[rgb]{0.81,0.36,0.00}{\textbf{#1}}}
\newcommand{\OtherTok}[1]{\textcolor[rgb]{0.56,0.35,0.01}{#1}}
\newcommand{\PreprocessorTok}[1]{\textcolor[rgb]{0.56,0.35,0.01}{\textit{#1}}}
\newcommand{\RegionMarkerTok}[1]{#1}
\newcommand{\SpecialCharTok}[1]{\textcolor[rgb]{0.81,0.36,0.00}{\textbf{#1}}}
\newcommand{\SpecialStringTok}[1]{\textcolor[rgb]{0.31,0.60,0.02}{#1}}
\newcommand{\StringTok}[1]{\textcolor[rgb]{0.31,0.60,0.02}{#1}}
\newcommand{\VariableTok}[1]{\textcolor[rgb]{0.00,0.00,0.00}{#1}}
\newcommand{\VerbatimStringTok}[1]{\textcolor[rgb]{0.31,0.60,0.02}{#1}}
\newcommand{\WarningTok}[1]{\textcolor[rgb]{0.56,0.35,0.01}{\textbf{\textit{#1}}}}
\usepackage{graphicx}
\makeatletter
\newsavebox\pandoc@box
\newcommand*\pandocbounded[1]{% scales image to fit in text height/width
  \sbox\pandoc@box{#1}%
  \Gscale@div\@tempa{\textheight}{\dimexpr\ht\pandoc@box+\dp\pandoc@box\relax}%
  \Gscale@div\@tempb{\linewidth}{\wd\pandoc@box}%
  \ifdim\@tempb\p@<\@tempa\p@\let\@tempa\@tempb\fi% select the smaller of both
  \ifdim\@tempa\p@<\p@\scalebox{\@tempa}{\usebox\pandoc@box}%
  \else\usebox{\pandoc@box}%
  \fi%
}
% Set default figure placement to htbp
\def\fps@figure{htbp}
\makeatother
\setlength{\emergencystretch}{3em} % prevent overfull lines
\providecommand{\tightlist}{%
  \setlength{\itemsep}{0pt}\setlength{\parskip}{0pt}}
\setcounter{secnumdepth}{-\maxdimen} % remove section numbering
\usepackage{bookmark}
\IfFileExists{xurl.sty}{\usepackage{xurl}}{} % add URL line breaks if available
\urlstyle{same}
\hypersetup{
  pdftitle={Homework 2},
  pdfauthor={Rowan Tichelaar},
  hidelinks,
  pdfcreator={LaTeX via pandoc}}

\title{Homework 2}
\author{Rowan Tichelaar}
\date{2026-02-23}

\begin{document}
\maketitle

\subsection{Homework 2}\label{homework-2}

\begin{Shaded}
\begin{Highlighting}[]
\NormalTok{ui }\OtherTok{\textless{}{-}} \FunctionTok{fluidPage}\NormalTok{(}
  \FunctionTok{titlePanel}\NormalTok{(}\StringTok{"Dynamic Scatter Plot of Mental Health Wellness of Gen Z"}\NormalTok{),}
  
  \FunctionTok{sidebarLayout}\NormalTok{(}
    \FunctionTok{sidebarPanel}\NormalTok{(}
      \CommentTok{\#Select input for x{-}axis}
      \FunctionTok{selectInput}\NormalTok{(}
        \AttributeTok{inputId =} \StringTok{"x\_var"}\NormalTok{,}
        \AttributeTok{label =} \StringTok{"Select X{-}axis variable:"}\NormalTok{,}
        \AttributeTok{choices =} \FunctionTok{names}\NormalTok{(data),}
        \AttributeTok{selected =} \FunctionTok{names}\NormalTok{(data)[}\DecValTok{7}\NormalTok{]}
\NormalTok{      ),}
      \CommentTok{\#Select input for y{-}axis}
      \FunctionTok{selectInput}\NormalTok{(}
        \AttributeTok{inputId =} \StringTok{"y\_var"}\NormalTok{,}
        \AttributeTok{label =} \StringTok{"Select Y{-}axis variable:"}\NormalTok{,}
        \AttributeTok{choices =} \FunctionTok{names}\NormalTok{(data),}
        \AttributeTok{selected =} \FunctionTok{names}\NormalTok{(data)[}\DecValTok{16}\NormalTok{]}
\NormalTok{      ),}
      \FunctionTok{sliderInput}\NormalTok{(}\StringTok{"alpha\_var"}\NormalTok{, }\StringTok{"Alpha"}\NormalTok{, }\AttributeTok{min =} \DecValTok{0}\NormalTok{, }\AttributeTok{max =} \DecValTok{1}\NormalTok{, }\AttributeTok{value =} \FloatTok{0.2}\NormalTok{),}
      
      \FunctionTok{checkboxInput}\NormalTok{(}
        \AttributeTok{inputId =} \StringTok{"show\_trend"}\NormalTok{, }
        \AttributeTok{label =} \StringTok{"Show Trend Line"}\NormalTok{, }
        \AttributeTok{value =} \ConstantTok{FALSE}
\NormalTok{      )}
\NormalTok{    ),}

    \CommentTok{\#Generate Plot}
    \FunctionTok{mainPanel}\NormalTok{(}
      \FunctionTok{plotOutput}\NormalTok{(}\StringTok{"scatter\_plot"}\NormalTok{)}
\NormalTok{    )}
\NormalTok{  )}
\NormalTok{)}

\NormalTok{server }\OtherTok{\textless{}{-}} \ControlFlowTok{function}\NormalTok{(input, output, session) \{}
  
  \CommentTok{\#Define the plot}
\NormalTok{  output}\SpecialCharTok{$}\NormalTok{scatter\_plot }\OtherTok{\textless{}{-}} \FunctionTok{renderPlot}\NormalTok{(\{}
    
\NormalTok{    p }\OtherTok{\textless{}{-}} \FunctionTok{ggplot}\NormalTok{(data, }\FunctionTok{aes}\NormalTok{(}\AttributeTok{x =}\NormalTok{ .data[[input}\SpecialCharTok{$}\NormalTok{x\_var]], }\AttributeTok{y =}\NormalTok{ .data[[input}\SpecialCharTok{$}\NormalTok{y\_var]])) }\SpecialCharTok{+} 
      \FunctionTok{geom\_point}\NormalTok{(}\AttributeTok{alpha =}\NormalTok{ input}\SpecialCharTok{$}\NormalTok{alpha\_var) }\SpecialCharTok{+}
      \FunctionTok{labs}\NormalTok{(}\AttributeTok{x =}\NormalTok{ input}\SpecialCharTok{$}\NormalTok{x\_var, }\AttributeTok{y =}\NormalTok{ input}\SpecialCharTok{$}\NormalTok{y\_var)}
    
    \CommentTok{\#Add smooth layer if show\_trend input is true}
    \ControlFlowTok{if}\NormalTok{ (input}\SpecialCharTok{$}\NormalTok{show\_trend) \{}
\NormalTok{      p }\OtherTok{\textless{}{-}}\NormalTok{ p }\SpecialCharTok{+} \FunctionTok{geom\_smooth}\NormalTok{(}\AttributeTok{method =} \StringTok{"lm"}\NormalTok{, }\AttributeTok{col =} \StringTok{"red"}\NormalTok{)}
\NormalTok{    \}}
    
\NormalTok{    p}
    
\NormalTok{  \})}
\NormalTok{\}}

\FunctionTok{shinyApp}\NormalTok{(ui, server)}
\end{Highlighting}
\end{Shaded}

I found this dataset on Kraggle, it was one of the trending data sets
and I found it interesting enough to use for this homework. I was trying
to find something interesting enough to represent for the app and I
couldn't decide so I decided to make it interactive in a way that
enabled you to decide what was being shown on the graph via x and y axis
input selectors.

I started with the input selectors for the x and y axis. It took a
little while and lots of Stack Overflow for me to figure out that I had
to make then strings, then I was finally able to get them into the
ggplot in the server section, as it had to be able to be entirely
reactive. Then I added a slider for the alpha, that was not nearly as
difficult. Then when I added the trend line and needed it to be able to
be disabled and enabled I had first experimented with another alpha
modifier but had decided that it was too complicated to set up a button
to output 1 and 0 instead of true and false and then get it to work in
the server. It also felt sorta of cheap and a cop out instead of making
a ``true'' disable/enable button. So I used the stuff we learned from
the faceting part of the lectures to figure out how to get it to work.
It's not true faceting and rather just adds it or not depending on the
button. Along with this, I also had to figure out a new way to handle
the x and y inputs to actual values, that once again had to be looked up
to figure it out.

This could easily be modified to include a dropdown to show other
graphs, whether or not you want it faceted and according to what, and
you could even add a text box that decides what data set to use to make
it a ``ez graph'' application. I also would like to clean it up a bit
more on the text department. It feels sorta like desmos except not just
for mathematical graphs and more for datasets.

The script can be found here:
\url{https://github.com/rowan-tich/Stat436-HW2/blob/main/HW2.Rmd} and
\url{https://uwmadison.box.com/s/r78oseme0zamlsjpk3fiwl3csbqvq2u3}

\end{document}
